% IEEE IoT-J: one paragraph, 150--250 words, no citations, no displayed math.
\begin{abstract}
Mobile underwater wireless sensor networks must keep link quality
current as vehicles move, yet a fixed neighbor-broadcast timer spends
scarce acoustic airtime while the channel is temporarily stable and
can still miss the instants at which a link-quality forecast has
diverged from the channel. This article presents Event-Triggered
Link-Quality and Topology Prediction (E-LQTP), an extension of the
link-quality-aware routing protocol LQARP. A dual-track extended
Kalman filter maintains a measurement-corrected relative-motion
estimate and a forecast of the same state that receives no
measurement update. Their innovation, gated by a persistence test and
a variance cap, emits a HELLO only when that forecast diverges from a
per-tick channel query. E-LQTP inherits LQARP's multipath channel
model, depth-progressive candidate set, and Bellhop integration; the
Q-learning reward is left unchanged at the validated operating point.
On the same Western Pacific site, with forty relays over six hundred
seconds and five seeds, every protocol reads the same live pose.
E-LQTP then matches LQARP's packet delivery, hops, delay, and energy
per data bit, while cutting control traffic to one-third of a
five-second HELLO grid under two-metre-per-second current drift and
to about one-half under a Rankine eddy. Additive predictive reward
terms do not improve the data plane. The control-plane gain is the
event-triggered HELLO under that shared pose, with the reward left
unchanged.
\end{abstract}
